% BANKS-SOURCE: pdf=258 print=248 kind=implicit-exercise id=I-APPC-001
\begin{exercise}[Hermiticity in the Weyl basis]{I-APPC-001}
Work in $D=4$ with the mostly-plus metric
$\eta=\operatorname{diag}(-1,+1,+1,+1)$ and the
Weyl matrices
\[
 \gamma^0=\begin{pmatrix}0&\mathbf 1_2\\ \mathbf 1_2&0\end{pmatrix},
 \qquad
\gamma^i=\begin{pmatrix}0&\sigma^i\\ -\sigma^i&0\end{pmatrix},
\qquad i=1,2,3,
\]
where the Pauli matrices obey $(\sigma^i)^\dagger=\sigma^i$ and
$\{\gamma^\mu,\gamma^\nu\}=-2\eta^{\mu\nu}$. Verify
directly, treating the time and space components separately, that
\[
 (\gamma^\mu)^\dagger=\gamma^0\gamma^\mu\gamma^0.
\]
Now let $U$ be unitary and define another representation by
$\gamma'^\mu=U^\dagger\gamma^\mu U$.  Prove that the primed matrices obey
the same adjoint relation with $\gamma'^0$, and identify the step that uses
unitarity.  Check along the way that $(\gamma^0)^2=\mathbf 1_4$ and
$\{\gamma^0,\gamma^i\}=0$.
\end{exercise}
